% GENERATED FILE. Edit canonical lessons, then run npm run generate:books.
% canonical-source-hash: fnv1a64:216feb485f06452b
% canonical-lessons: PT-C512-pavilhao, PT-C512-basquetebol, PT-C512-andebol, PT-C512-voleibol, PT-C512-ginastica

\chapter{No pavilhão}
\label{ch:pt-no-pavilhao}
\hlchaptermodality{512}

\begin{chapteropening}
\textbf{By the end of this chapter:} \emph{I can talk about indoor sports.}

Complete the last lesson of chapter 512: I can talk about indoor sports.
\end{chapteropening}

\section[o pavilhão]{o pavilhão — sports hall; pavilion}

\label{lesson:PT-C512-pavilhao}



\begin{quote}
Before the new one: say the Portuguese for podium, then the Portuguese for record.
\end{quote}

\subsection*{You'll want to know: o pavilhão}
\textbf{o pavilhão} — \textquotedblleft{}sports hall; pavilion\textquotedblright{}.

\begin{quote}
Os treinos são no pavilhão da escola.
\end{quote}

\begin{quote}
\emph{Training is in the school's sports hall.}
\end{quote}

\emph{Pavilhão desportivo} — sports hall.

\begin{cousinweb}
\emph{pavilhão} was borrowed from French \emph{pavillon} (EtymDB), from Latin \emph{pāpiliō} \textquotedblleft{}a butterfly\textquotedblright{}, and so \textquotedblleft{}a tent\textquotedblright{}, whose spread cloth looked like wings.
\end{cousinweb}

\begin{grammarlens}[title={What you've built}]
One more word for telling, explaining and describing at B1.
\end{grammarlens}

\subsection*{Your turn}
\begin{itemize}
\raggedright
  \item \emph{Say it:} \emph{o pavilhão}
  \item \emph{Say it:} \emph{o pavilhão}, once more
  \item \emph{Say it:} a sentence of your own with \emph{o pavilhão}
  \item \emph{Recall:} say the Portuguese for podium, then the Portuguese for record, then say \emph{o pavilhão} again
\end{itemize}

\begin{tcolorbox}[breakable,colback=teal!4,colframe=teal!35!black,title={Before you move on}]
What does o pavilhão mean? (\textquotedblleft{}Sports hall; pavilion\textquotedblright{}.) Say it once more.
\end{tcolorbox}
\section[o basquetebol]{o basquetebol — basketball}

\label{lesson:PT-C512-basquetebol}



\begin{quote}
Before the new one: say the Portuguese for record, then the Portuguese for sports hall.
\end{quote}

\subsection*{You'll want to know: o basquetebol}
\textbf{o basquetebol} — \textquotedblleft{}basketball\textquotedblright{}.

\begin{quote}
Ele é alto e joga basquetebol.
\end{quote}

\begin{quote}
\emph{He is tall and plays basketball.}
\end{quote}

In Brazil: \emph{o basquete}.

\begin{cousinweb}
\emph{basquetebol} was borrowed from English \emph{basketball} (EtymDB).
\end{cousinweb}

\begin{grammarlens}[title={What you've built}]
One more word for telling, explaining and describing at B1.
\end{grammarlens}

\subsection*{Your turn}
\begin{itemize}
\raggedright
  \item \emph{Say it:} \emph{o basquetebol}
  \item \emph{Say it:} \emph{o basquetebol}, once more
  \item \emph{Say it:} a sentence of your own with \emph{o basquetebol}
  \item \emph{Recall:} say the Portuguese for record, then the Portuguese for sports hall, then say \emph{o basquetebol} again
\end{itemize}

\begin{tcolorbox}[breakable,colback=teal!4,colframe=teal!35!black,title={Before you move on}]
What does o basquetebol mean? (\textquotedblleft{}Basketball\textquotedblright{}.) Say it once more.
\end{tcolorbox}
\section[o andebol]{o andebol — handball}

\label{lesson:PT-C512-andebol}



\begin{quote}
Before the new one: say the Portuguese for sports hall, then the Portuguese for basketball.
\end{quote}

\subsection*{You'll want to know: o andebol}
\textbf{o andebol} — \textquotedblleft{}handball\textquotedblright{}.

\begin{quote}
Joga andebol no clube do bairro.
\end{quote}

\begin{quote}
\emph{He plays handball at the local club.}
\end{quote}

In Brazil: \emph{o handebol}.

\begin{cousinweb}
\emph{andebol} was borrowed from English \emph{handball} (EtymDB), adapted to Portuguese spelling.
\end{cousinweb}

\begin{grammarlens}[title={What you've built}]
One more word for telling, explaining and describing at B1.
\end{grammarlens}

\subsection*{Your turn}
\begin{itemize}
\raggedright
  \item \emph{Say it:} \emph{o andebol}
  \item \emph{Say it:} \emph{o andebol}, once more
  \item \emph{Say it:} a sentence of your own with \emph{o andebol}
  \item \emph{Recall:} say the Portuguese for sports hall, then the Portuguese for basketball, then say \emph{o andebol} again
\end{itemize}

\begin{tcolorbox}[breakable,colback=teal!4,colframe=teal!35!black,title={Before you move on}]
What does o andebol mean? (\textquotedblleft{}Handball\textquotedblright{}.) Say it once more.
\end{tcolorbox}
\section[o voleibol]{o voleibol — volleyball}

\label{lesson:PT-C512-voleibol}



\begin{quote}
Before the new one: say the Portuguese for basketball, then the Portuguese for handball.
\end{quote}

\subsection*{You'll want to know: o voleibol}
\textbf{o voleibol} — \textquotedblleft{}volleyball\textquotedblright{}.

\begin{quote}
Na praia jogamos voleibol.
\end{quote}

\begin{quote}
\emph{On the beach we play volleyball.}
\end{quote}

Also \emph{o vólei}.

\begin{cousinweb}
\emph{voleibol} was borrowed from English \emph{volleyball} (EtymDB).
\end{cousinweb}

\begin{grammarlens}[title={What you've built}]
One more word for telling, explaining and describing at B1.
\end{grammarlens}

\subsection*{Your turn}
\begin{itemize}
\raggedright
  \item \emph{Say it:} \emph{o voleibol}
  \item \emph{Say it:} \emph{o voleibol}, once more
  \item \emph{Say it:} a sentence of your own with \emph{o voleibol}
  \item \emph{Recall:} say the Portuguese for basketball, then the Portuguese for handball, then say \emph{o voleibol} again
\end{itemize}

\begin{tcolorbox}[breakable,colback=teal!4,colframe=teal!35!black,title={Before you move on}]
What does o voleibol mean? (\textquotedblleft{}Volleyball\textquotedblright{}.) Say it once more.
\end{tcolorbox}
\section[a ginástica]{a ginástica — gymnastics; exercise}

\label{lesson:PT-C512-ginastica}



\begin{quote}
Before the new one: say the Portuguese for handball, then the Portuguese for volleyball.
\end{quote}

\subsection*{You'll want to know: a ginástica}
\textbf{a ginástica} — \textquotedblleft{}gymnastics; exercise\textquotedblright{}.

\begin{quote}
Faço ginástica três vezes por semana.
\end{quote}

\begin{quote}
\emph{I exercise three times a week.}
\end{quote}

\emph{Fazer ginástica} — to do exercise.

\begin{cousinweb}
\emph{ginástica} is a learned borrowing through Latin \emph{gymnasticus} \textquotedblleft{}of bodily exercise\textquotedblright{} (Lewis \& Short) from Greek \emph{gymnastikós}, built on \emph{gymnós} \textquotedblleft{}naked\textquotedblright{}: Greek athletes trained unclothed.
\end{cousinweb}

\begin{grammarlens}[title={What you've built}]
One more word for telling, explaining and describing at B1.
\end{grammarlens}

\subsection*{Your turn}
\begin{itemize}
\raggedright
  \item \emph{Say it:} \emph{a ginástica}
  \item \emph{Say it:} \emph{a ginástica}, once more
  \item \emph{Say it:} a sentence of your own with \emph{a ginástica}
  \item \emph{Recall:} say the Portuguese for handball, then the Portuguese for volleyball, then say \emph{a ginástica} again
\end{itemize}

\begin{tcolorbox}[breakable,colback=teal!4,colframe=teal!35!black,title={Before you move on}]
What does a ginástica mean? (\textquotedblleft{}Gymnastics; exercise\textquotedblright{}.) Say it once more.
\end{tcolorbox}
